% page 521 of the facsimile (image p-520 of the scan)
% block 170.2 x 237.7 mm ; left margin 31.8 mm
\pgc{10.160mm}{0.000mm}{
\l{\cel{57}513}
\l{}
\l{\sou{sarki}no. (\sou{mikrobio}l.) Bakterio saprofita quan onu renkontras en}
\l{\cel{3}la gangreno pulmonala ed ula morbi di la stomako. -}
\l{}
\l{\sou{sarklar}. (\sou{trans.}) Netigar (sulo-peco) per extirpar la herbi mala.}
\l{\cel{3}- FI.}
\l{}
\l{\sou{sarkofago}.I. Sarko de petro en qua la antiqui depozis la korpi}
\l{\cel{3}mortinta. - II. Parto di monumento funerala qua simulas petro}
\l{\cel{3}sarko qua kontenas la kadavro. - DEFIRS.}
\l{}
\l{\sou{sarkomo}. (\sou{patol.}) Surkreskajo de karno. - DEFIS.}
\l{}
\l{\sou{sarkoplasmo.(histol.}) Protoplasmo granuloza di la fibro muskulala,}
\l{\cel{3}qua cirkondas omna nukleo, e su extensas, per traco-linei plu}
\l{\cel{3}o min dina e longa, inter la grupi quin formacas la fibreti}
\l{\cel{3}muskulala. -}
\l{}
\l{\sou{sarmento.}\cel{2}(\sou{bot.}) Ligno quan la vito produktas singlayare. - EFIS.}
\l{}
\l{\sou{sarsaparelo}. (\sou{bot.}) Planto di qua la radiko esas purigiva medi-}
\l{\cel{3}cinale. - L.\cel{1}\sou{smilax}. DEFIRS.}
\l{}
\l{\sou{sat (e)}. Tam (multa) kam bezonata.- L.}
\l{}
\l{\sou{satano}. (\sou{kristan.}) La nomo di la diablo. - DEFIRS.}
\l{}
\l{\sou{satelito}. I. Viro armoza, dop chefo, por exekutar lua imperi. -}
\l{\cel{3}(\sou{anke metaf}.). II. (\sou{astron.}) Planeto acesora qua jiras cirkum}
\l{\cel{3}planeto e sequas ica en lua rotaco. - DEFIRS.}
\l{}
\l{\sou{satino}. Stofo ek silko, plana, en qua la wefto ne esas videbla}
\l{\cel{3}sur la averso, e qua esas briligita per la preso da cilindro.}
\l{\cel{3}- DEFIR.}
\l{}
\l{\sou{satiro}. I. (\sou{literaturo Latina}). Verko ek prozo e versi, en qua}
\l{\cel{3}la autoro blamegas la mori di publiko (+\sou{publico}).- II. Verko}
\l{\cel{3}ek versi ube la poeto atakas la vicii e la mokindaji di lua sa-}
\l{\cel{3}mepokani. - DEFIRS.}
\l{}
\l{\sou{satiruso}. I. (\sou{mitol.}) Mi-deo qua habitas la boski, kun korpo longa}
\l{\cel{3}pila, korni, gambi e pedi kaprulala. - II. (\sou{metaf.}) Viro laciva}
\l{\cel{3}debochema. - DEFIRS.}
\l{}
\l{\sou{satisfacar}. (\sou{trans., per}) I. Igar (ulu) en stando agreabla per}
\l{\cel{3}agar o facar por lu to quon lu expektas o deziras. - II.(\sou{ulo})}
\l{\cel{3}Realigar to quon postulas ulo. - DEFIS.}
\l{}
\l{\sou{\textquotedbl{}satisfecit\textquotedbl{}}. (\sou{vorto skolanala)}.\cel{2}Paperfolio per qua docisto}
\l{\cel{3}atestas satisfacesir da la agi o faci da lua docario. -}
\l{}
\l{\sou{satrapo.}\cel{2}I. (\sou{epoki antiqua}).Guvernisto qua darfis agar suverene}
\l{\cel{3}en provinco ye la nomo di la rejulo di Persia. - II. (nia-epoke)}
\l{\cel{3}Persono qua, en sua ofico-resortiso, agas despote. - DEFIRS.}
}
